% =====================================================================
%  7. Marco teórico
% =====================================================================
\section{Marco teórico}
\label{sec:marco-teorico}

El marco teórico no describe las tecnologías empleadas --- eso corresponde a la
Parte~IV --- sino los \textbf{constructos} con los que esta investigación piensa el
problema. Cada uno se presenta con tres elementos: qué afirma, de dónde proviene y
qué papel cumple en el diseño de la investigación.

\subsection{Calidad del corpus como propiedad relacional}
\label{sec:teo-calidad}

La literatura de calidad de recursos educativos trata la calidad como una propiedad
\emph{del recurso}: LORI evalúa nueve dimensiones de un objeto de aprendizaje
\parencite{leacock2007framework} y el modelo por capas de
\textcite{tabares2017modelo} integra tres visiones sobre un mismo objeto. Esta
investigación introduce un desplazamiento: sobre el corpus de un curso, parte de la
calidad no reside en ningún documento sino en la \textbf{relación entre documentos}.

Dos documentos individualmente exactos pueden ser conjuntamente inconsistentes. Un
documento sin errores puede estar desactualizado respecto de otro que lo sustituyó
sin retirarlo. Un contenido correcto puede ser redundante respecto de otro
equivalente. En los tres casos, la unidad evaluable es el par o el conjunto, no el
elemento.

\begin{estado}
\small
\textbf{Papel en la investigación.} Este desplazamiento determina la unidad de
análisis del benchmark: las instancias etiquetadas no son documentos sino
\emph{relaciones} entre fragmentos, con su evidencia asociada. Determina también por
qué los instrumentos de la línea A no pueden adoptarse directamente como medida de
resultado.
\end{estado}

\subsection{La consistencia como relación de inferencia de tres valores}
\label{sec:teo-inferencia}

El constructo que permite formalizar la detección proviene de la inferencia en
lenguaje natural. Dada una afirmación y un texto de referencia, la relación admite
tres valores: el texto \textbf{implica} la afirmación, la \textbf{contradice}, o
\textbf{no aporta información suficiente} para decidir. La tarea de
\textcite{thorne2018fever} fija esas tres clases con un criterio de anotación
conservador --- sin evidencia explícita, la clase es la tercera, aunque el anotador
crea saber la respuesta --- y \textcite{koreeda2021contractnli} la trasladan al nivel
documental exigiendo que se identifique la evidencia que sustenta la decisión.

La adopción de este constructo tiene tres consecuencias:

\begin{enumerate}
  \item \textbf{La clase neutral no es un residuo.} \enquote{No encontré con qué
        comparar esto} es una salida informativa y distinta de \enquote{esto no se
        contradice con nada}. Colapsarlas produce alertas que el docente no puede
        interpretar.
  \item \textbf{La verificación es relativa a la fuente.} Lo que se decide no es si
        una afirmación es verdadera, sino si está respaldada por el corpus
        consultado. \textcite{thorne2018fever} evitan deliberadamente el vocabulario
        de verdadero y falso por esta razón.
  \item \textbf{La dificultad es asimétrica entre clases.} La contradicción es la
        clase más difícil \parencite{koreeda2021contractnli} y la neutral la peor
        reconocida \parencite{gubelmann2024capturing}. Un desempeño agregado oculta
        esta asimetría, por lo que las métricas deben reportarse por clase.
\end{enumerate}

\subsection{Anclaje en evidencia}
\label{sec:teo-evidencia}

El anclaje --- que cada afirmación del sistema remita a fragmentos localizables del
corpus --- cumple aquí tres funciones que conviene no confundir.

Como \textbf{mecanismo de mitigación}, reduce la dependencia de la memoria
paramétrica del modelo \parencite{lewis2020retrieval}, aunque no elimina la
alucinación \parencite{huang2025survey}. Como \textbf{condición de verificabilidad},
permite que el usuario localice la referencia original y compruebe la afirmación,
manteniendo observable un proceso que de otro modo sería opaco
\parencite[p.~14]{gao2024retrieval}. Como \textbf{objeto de medición}, convierte la
fidelidad en una propiedad evaluable: puede verificarse, afirmación por afirmación,
si la evidencia citada la sustenta \parencite{min2023factscore,es2024ragas}.

\begin{advertencia}
La tercera función depende de la primera pero no se sigue de ella. Que un sistema
exhiba evidencia no implica que la evidencia sustente lo que se afirma. Un sistema
puede citar fragmentos reales y aun así derivar de ellos una conclusión que no se
sigue. Por eso la fidelidad es una variable a medir y no una propiedad garantizada
por el diseño.
\end{advertencia}

\subsection{Confianza, calibración y confianza apropiada}
\label{sec:teo-confianza}

Tres constructos distintos se confunden con frecuencia bajo la palabra
\enquote{confianza}. Esta investigación los mantiene separados.

\begin{description}[style=nextline,leftmargin=0.8cm]
  \item[Similitud] Medida de cercanía entre representaciones vectoriales de dos
        textos \parencite{reimers2019sentencebert}. Indica proximidad temática. No es
        una probabilidad ni indica corrección.
  \item[Confianza estimada] Puntaje que el sistema asigna a su propia salida. Puede
        construirse de muchas maneras, incluidas combinaciones heurísticas de
        señales. Un puntaje normalizado al intervalo $[0,1]$ no es por ello una
        probabilidad.
  \item[Calibración] Correspondencia estadística entre la confianza estimada y la
        frecuencia observada de acierto. Es una propiedad \emph{demostrable} y
        \emph{medible}, no una característica de diseño: exige comparar confianza y
        corrección sobre un conjunto independiente con etiquetas verdaderas
        \parencite{guo2017calibration}.
  \item[Confianza apropiada] Propiedad de la relación entre la persona y el sistema:
        que el usuario confíe cuando el sistema acierta y desconfíe cuando falla
        \parencite{mehrotra2024systematic}. No equivale a aceptación ni a
        dependencia.
\end{description}

La relación entre los cuatro constructos no es de grado sino de naturaleza. El
hallazgo de \textcite{guo2017calibration} de que exactitud y calibración no se
optimizan con los mismos parámetros implica que un sistema puede mejorar su detección
y empeorar simultáneamente la calidad de sus señales de confianza. Y el de
\textcite{mehrotra2024systematic}, que un puntaje mal calibrado puede inducir
confianza inapropiada, implica que exhibir un puntaje no es neutro: puede empeorar la
decisión del docente.

\subsection{Retroalimentación humana como adaptación contextual}
\label{sec:teo-feedback}

La incorporación de decisiones previas del docente al comportamiento futuro del
sistema se apoya en la estrategia de \textbf{memoria de retroalimentación} descrita
por \textcite{fernandes2023bridging}: almacenar las decisiones de sesiones anteriores
y, ante una entrada nueva, recuperar las correspondientes a entradas similares para
condicionar la generación. Es una adaptación en tiempo de inferencia: no modifica los
parámetros del modelo.

Esta precisión terminológica importa. Lo que el sistema hace \textbf{no es
aprendizaje} en el sentido de ajuste de parámetros, y describirlo como tal sería
incorrecto. Es recuperación de precedentes. Los riesgos que \textcite{fernandes2023bridging}
documentan para la recolección de retroalimentación --- desacuerdo entre expertos,
sesgo de anclaje, sesgo de positividad, retroalimentación contradictoria --- se
aplican íntegramente a un repositorio de decisiones docentes y constituyen amenazas a
la validez del experimento correspondiente.

\subsection{Procedencia y delegación de autoridad}
\label{sec:teo-procedencia}

El modelo de procedencia del \textcite{moreau2013provdm} aporta dos elementos.

El primero es una \textbf{restricción sobre lo que un registro permite afirmar}: la
derivación no se infiere de la coocurrencia. Que una actividad haya usado una entidad
y generado otra no establece que la segunda derive de la primera. Trasladado a este
proyecto: registrar que el sistema consultó unos fragmentos y produjo una sugerencia
no demuestra que la sugerencia se apoye en ellos. La trazabilidad que puede afirmarse
es la del proceso, no la del fundamento; esta última debe medirse.

El segundo es la noción de \textbf{delegación}: un agente actúa en nombre de otro que
conserva la responsabilidad por el resultado. Formaliza con precisión la relación que
este proyecto sostiene entre el sistema y el docente, y da forma registrable a la
cadena propuesta $\rightarrow$ evidencia $\rightarrow$ decisión $\rightarrow$
responsable.

\subsection{Agencia humana como restricción, no como característica}
\label{sec:teo-agencia}

La posición de que la decisión permanece en el docente no es una preferencia de
diseño de este proyecto: es un requisito derivado de dos marcos normativos. El
\textcite{nist2023airmf} exige documentar los procesos de supervisión humana y los
límites de conocimiento del sistema, y recomienda recoger datos sobre la frecuencia y
las razones con que las personas anulan las salidas de la IA. La guía de
\textcite{miao2023guidance} establece que no debe cederse la responsabilidad humana
en decisiones de alto impacto y que el diseñador humano debe verificar el
conocimiento factual de lo que la IA sugiera.

La consecuencia metodológica es que los motivos de rechazo dejan de ser un registro
operativo del sistema y se convierten en \textbf{dato de investigación}.

\subsection{Design Science Research como teoría del método}
\label{sec:teo-dsr}

El marco metodológico se toma de \textcite{hevner2004design} y
\textcite{peffers2007design}. Del primero, la exigencia de que la contribución sea
identificable respecto de la base de conocimiento y el repertorio de métodos de
evaluación del artefacto. Del segundo, la admisión explícita de puntos de entrada
distintos al ciclo, entre ellos el centrado en diseño y desarrollo, y el carácter
iterativo y no rígido del proceso.

La advertencia de \textcite{hevner2004design} sobre el carácter perecedero de la
investigación en diseño --- los avances tecnológicos pueden dejarla obsoleta con
rapidez --- se traduce en una decisión concreta de este proyecto: las contribuciones
que se persiguen son el \emph{protocolo}, el \emph{benchmark} anotado y el
\emph{conocimiento sobre qué componentes explican el comportamiento}, artefactos
menos perecederos que una implementación concreta atada a versiones de modelos.
